\documentclass[11pt,a4paper]{article}
% ══════════════════════════════════════════════════════════════════════════════
% Edexcel Further Mathematics (Core Pure 1) --- Shared LaTeX Preamble
% Elegant, Minimalist & Aesthetic Academic Design System
% ══════════════════════════════════════════════════════════════════════════════

\usepackage[a4paper, top=22mm, bottom=22mm, left=20mm, right=20mm]{geometry}
\usepackage{amsmath, amssymb, amsthm}
\usepackage{enumitem}
\usepackage{booktabs}
\usepackage{array}
\usepackage{parskip}
\usepackage{microtype}
\usepackage{fancyhdr}
\usepackage{titlesec}
\usepackage{mdframed}
\usepackage{xcolor}
\usepackage[hidelinks]{hyperref}
\usepackage{graphicx}
\usepackage{tikz}
\usepackage{pgfplots}
\pgfplotsset{compat=1.18}

% ── Refined Palette ───────────────────────────────────────────────────────────
\definecolor{primary}{gray}{0.10}
\definecolor{darkgray}{gray}{0.20}
\definecolor{medgray}{gray}{0.50}
\definecolor{lightgray}{gray}{0.96}
\definecolor{boxgray}{gray}{0.85}
\definecolor{accentline}{gray}{0.75}

% ── Section formatting ──────────────────────────────────────────────────────────
\titleformat{\section}[block]
  {\normalfont\bfseries\color{primary}}{}
  {0pt}{}[\vspace{1pt}{\color{accentline}\rule{\textwidth}{0.4pt}}\vspace{4pt}]
\titlespacing*{\section}{0pt}{14pt}{6pt}

% ── List formatting ─────────────────────────────────────────────────────────────
\setlist[enumerate,1]{label=\textbf{\arabic*.}, leftmargin=1.8em, itemsep=14pt}
\setlist[enumerate,2]{label=\textbf{(\alph*)}, leftmargin=1.8em, itemsep=6pt}

% ── Header / footer (Overlaps permanently eliminated) ─────────────────────────
\pagestyle{fancy}
\newcommand{\CPHeader}[2]{%
  \fancyhf{}%
  \setlength{\headheight}{22pt}%
  \renewcommand{\headrulewidth}{0.3pt}%
  \renewcommand{\headrule}{\hbox to\headwidth{\color{accentline}\leaders\hrule height \headrulewidth\hfill}}%
  \fancyhead[L]{\footnotesize\scshape\color{medgray} Core Pure 1: #1}%
  \fancyhead[R]{\footnotesize\color{medgray} Chapter \##2}%
  \renewcommand{\footrulewidth}{0.3pt}%
  \renewcommand{\footrule}{\hbox to\headwidth{\color{accentline}\leaders\hrule height \footrulewidth\hfill}}%
  \fancyfoot[L]{\footnotesize\color{medgray} Pearson Edexcel A-Level Further Mathematics $\cdot$ Core Pure 1}%
  \fancyfoot[C]{}% No center text to ensure zero collision with long titles
  \fancyfoot[R]{\footnotesize\color{medgray} \textbf{Henry.K} \enspace$\vert$\enspace Page \thepage}%
}

% ── Elegant Aesthetic Title Block ──────────────────────────────────────────────
\newcommand{\CPTitleBlock}[2]{%
  \vspace*{-8pt}
  \begin{center}
    {\footnotesize\scshape\color{medgray} Pearson Edexcel A-Level Further Mathematics \enspace$\cdot$\enspace Core Pure 1}\\[4pt]
    {\LARGE\bfseries\color{primary} #1}\\[4pt]
    {\small\color{darkgray} #2}\\[6pt]
    {\footnotesize\color{medgray} Problem Series authored by \textbf{Henry.K} \quad$\cdot$\quad 10 Exam Archetypes \quad$\cdot$\quad 65 Total Marks}\\[8pt]
    {\color{accentline}\rule{\textwidth}{0.5pt}}
  \end{center}
  \vspace{-6pt}
  \noindent{\footnotesize\itshape \textbf{Instructions:} Answer all ten questions. Full algebraic working must be shown for method marks; unsupported answers from calculators will not gain full credit. Give exact values in surd or fractional form unless specified otherwise.}\par
  \vspace{3pt}
  {\color{accentline}\rule{\textwidth}{0.5pt}}
  \vspace{8pt}
}

% ── Answer / Method / Investigate-further boxes ────────────────────────────────
\newmdenv[
  backgroundcolor=lightgray,
  linecolor=boxgray,
  linewidth=0.5pt,
  roundcorner=3pt,
  leftmargin=0pt, rightmargin=0pt,
  innerleftmargin=10pt, innerrightmargin=10pt,
  innertopmargin=6pt, innerbottommargin=6pt,
  skipabove=5pt, skipbelow=5pt
]{ansbox}

\newmdenv[
  backgroundcolor=lightgray!40,
  linecolor=darkgray,
  linewidth=1.2pt,
  leftline=true, rightline=false, topline=false, bottomline=false,
  leftmargin=0pt, rightmargin=0pt,
  innerleftmargin=10pt, innerrightmargin=8pt,
  innertopmargin=5pt, innerbottommargin=5pt,
  skipabove=5pt, skipbelow=7pt
]{invbox}

\newcommand{\ans}[1]{%
  \begin{ansbox}
    \textbf{Answer:} #1
  \end{ansbox}}

\newcommand{\method}[1]{%
  {\small\textbf{Method Criteria:}\par\vspace{2pt}#1}\par}

\newcommand{\inv}[1]{%
  \begin{invbox}
    \textbf{\small Deep Investigation:} {\small #1}
  \end{invbox}}

\newcommand{\qn}[1]{\textbf{#1.}\;}

\usetikzlibrary{arrows.meta,calc}
\CPHeader{Matrices \& Linear Systems --- Solutions}{6}

\begin{document}
\thispagestyle{plain}

\CPTitleBlock{Matrices \& Linear Systems}{Solutions \& Mark Scheme \#6 \enspace$\cdot$\enspace Determinants, Inverses, Adjugates, Consistency \& 3D Planes}

\section*{Section A \quad Procedural \& Direct Foundations}

%── Q1 [DIR] ─────────────────────────────────────────────────────────────
\qn{1} \textbf{Determinant Calculation and Multiplicative Properties}
\ans{(a) $\det(\mathbf{A}) = 10$; \quad (b)(i) $\det(\mathbf{B}^{-1}) = \frac{1}{12}$; \quad (ii) $\det(2\mathbf{B}) = 96$.}

\method{
  \textbf{(a)} Expand along Row 1:
  \[
  \det(\mathbf{A}) = 2\begin{vmatrix} 3 & 2 \\ 2 & 1 \end{vmatrix} - 1\begin{vmatrix} 1 & 2 \\ 3 & 1 \end{vmatrix} + (-1)\begin{vmatrix} 1 & 3 \\ 3 & 2 \end{vmatrix}
  \]
  \begin{itemize}[leftmargin=1.5em, itemsep=2pt]
    \item \textbf{M1}: Correct cofactor expansion along any row or column.
    \item \textbf{A1}: $2(3 - 4) - 1(1 - 6) - 1(2 - 9) = -2 + 5 + 7 = 10$.
  \end{itemize}
  \textbf{(b)(i)}
  \begin{itemize}[leftmargin=1.5em, itemsep=2pt]
    \item \textbf{M1}: Apply $\det(\mathbf{A}\mathbf{B}) = \det(\mathbf{A})\det(\mathbf{B}) \implies 10\det(\mathbf{B}) = 120 \implies \det(\mathbf{B}) = 12$.
    \item \textbf{A1}: Use $\det(\mathbf{B}^{-1}) = \frac{1}{\det(\mathbf{B})} = \frac{1}{12}$.
  \end{itemize}
  \textbf{(b)(ii)}
  \begin{itemize}[leftmargin=1.5em, itemsep=2pt]
    \item \textbf{B1}: For a $3 \times 3$ matrix, $\det(k\mathbf{B}) = k^3 \det(\mathbf{B})$, so $\det(2\mathbf{B}) = 2^3 \times 12 = 8 \times 12 = 96$.
  \end{itemize}
}

\inv{Why does $\det(k\mathbf{B}) = k^n \det(\mathbf{B})$ for an $n \times n$ matrix? In geometric terms, multiplying every entry of an $n \times n$ matrix by $k$ scales all $n$ basis vectors by $k$, so the hypervolume scales by $k^n$. What happens to the determinant if only a single row or column is scaled by $k$?}

\vspace{10pt}

%── Q2 [DIR] ─────────────────────────────────────────────────────────────
\qn{2} \textbf{Inversion of a $3\times3$ Matrix via Adjugate}
\ans{(a) $\det(\mathbf{M}) = -4$; \quad (b) $\mathbf{C} = \begin{pmatrix} 1 & -3 & 1 \\ -3 & 1 & 1 \\ 1 & 1 & -3 \end{pmatrix}$; \quad (c) $\mathbf{M}^{-1} = \frac{1}{4}\begin{pmatrix} -1 & 3 & -1 \\ 3 & -1 & -1 \\ -1 & -1 & 3 \end{pmatrix}$.}

\method{
  \textbf{(a)}
  \begin{itemize}[leftmargin=1.5em, itemsep=2pt]
    \item \textbf{B1}: $\det(\mathbf{M}) = 1(2 - 1) - 2(4 - 1) + 1(2 - 1) = 1 - 6 + 1 = -4$.
  \end{itemize}
  \textbf{(b)}
  \begin{itemize}[leftmargin=1.5em, itemsep=2pt]
    \item \textbf{M1}: Compute minors for each position:
    $M_{11}=1, M_{12}=3, M_{13}=1$; $M_{21}=3, M_{22}=1, M_{23}=-1$; $M_{31}=1, M_{32}=-1, M_{33}=-3$.
    \item \textbf{A1}: Correct sign checkerboard pattern $C_{ij} = (-1)^{i+j} M_{ij}$.
    \item \textbf{A1}: Complete cofactor matrix $\mathbf{C} = \begin{pmatrix} 1 & -3 & 1 \\ -3 & 1 & 1 \\ 1 & 1 & -3 \end{pmatrix}$.
  \end{itemize}
  \textbf{(c)}
  \begin{itemize}[leftmargin=1.5em, itemsep=2pt]
    \item \textbf{M1}: Transpose cofactors to get $\operatorname{adj}(\mathbf{M}) = \mathbf{C}^T$ and apply $\mathbf{M}^{-1} = \frac{1}{\det(\mathbf{M})}\operatorname{adj}(\mathbf{M})$.
    \item \textbf{A1}: Since $\mathbf{C}$ is symmetric, $\operatorname{adj}(\mathbf{M}) = \mathbf{C}$. Hence $\mathbf{M}^{-1} = -\frac{1}{4}\begin{pmatrix} 1 & -3 & 1 \\ -3 & 1 & 1 \\ 1 & 1 & -3 \end{pmatrix} = \frac{1}{4}\begin{pmatrix} -1 & 3 & -1 \\ 3 & -1 & -1 \\ -1 & -1 & 3 \end{pmatrix}$.
  \end{itemize}
}

\inv{Notice that $\mathbf{M}$ is a symmetric matrix ($\mathbf{M} = \mathbf{M}^T$). Does the inverse of any invertible symmetric matrix always remain symmetric? Prove this algebraically using the properties of transposes: $(\mathbf{M}^{-1})^T = (\mathbf{M}^T)^{-1}$.}

\vspace{10pt}

%── Q3 [DIR] ─────────────────────────────────────────────────────────────
\qn{3} \textbf{Solving Linear Systems using Matrix Inverses}
\ans{(a) $\begin{pmatrix} 1 & 2 & 1 \\ 2 & 1 & 1 \\ 1 & 1 & 2 \end{pmatrix}\begin{pmatrix} x \\ y \\ z \end{pmatrix} = \begin{pmatrix} 8 \\ 7 \\ 9 \end{pmatrix}$; \quad (b) $x = 1,\; y = 2,\; z = 3$.}

\method{
  \textbf{(a)}
  \begin{itemize}[leftmargin=1.5em, itemsep=2pt]
    \item \textbf{B1}: State correct matrix equation $\mathbf{M}\mathbf{x} = \mathbf{b}$.
  \end{itemize}
  \textbf{(b)}
  \begin{itemize}[leftmargin=1.5em, itemsep=2pt]
    \item \textbf{M1}: Premultiply by $\mathbf{M}^{-1}$: $\mathbf{x} = \mathbf{M}^{-1}\mathbf{b}$.
    \item \textbf{M1}: Execute matrix multiplication:
    \[
    \begin{pmatrix} x \\ y \\ z \end{pmatrix} = \frac{1}{4} \begin{pmatrix} -1 & 3 & -1 \\ 3 & -1 & -1 \\ -1 & -1 & 3 \end{pmatrix} \begin{pmatrix} 8 \\ 7 \\ 9 \end{pmatrix} = \frac{1}{4} \begin{pmatrix} -8 + 21 - 9 \\ 24 - 7 - 9 \\ -8 - 7 + 27 \end{pmatrix}
    \]
    \item \textbf{A1}: Vector product yields $\frac{1}{4}\begin{pmatrix} 4 \\ 8 \\ 12 \end{pmatrix}$.
    \item \textbf{A1}: $x = 1,\; y = 2,\; z = 3$.
  \end{itemize}
}

\inv{Why is finding the matrix inverse inefficient computationally for large $n \times n$ systems compared to Gaussian elimination or LU decomposition? Contrast the $\mathcal{O}(n^3)$ operation counts.}

\newpage
\section*{Section B \quad Reverse Parameters \& Analysis}

%── Q4 [REV] ─────────────────────────────────────────────────────────────
\qn{4} \textbf{Singular Matrices and Existence of Inverse}
\ans{(a) Proof complete; \quad (b) $k = 0$ or $k = 3$; \quad (c) $\{k \in \mathbb{R} : k \neq 0 \text{ and } k \neq 3\}$.}

\method{
  \textbf{(a)}
  \begin{itemize}[leftmargin=1.5em, itemsep=2pt]
    \item \textbf{M1}: Expand determinant: $1(k - 4) - 2(2 - 4) + k(2 - k)$.
    \item \textbf{A1}: Simplify: $(k - 4) - 2(-2) + (2k - k^2) = k - 4 + 4 + 2k - k^2$.
    \item \textbf{A1}: Combine terms to obtain $3k - k^2 = k(3 - k)$.
  \end{itemize}
  \textbf{(b)}
  \begin{itemize}[leftmargin=1.5em, itemsep=2pt]
    \item \textbf{B1}: $\mathbf{P}$ is singular $\iff \det(\mathbf{P}) = 0 \implies k(3 - k) = 0 \implies k = 0$ or $k = 3$.
  \end{itemize}
  \textbf{(c)}
  \begin{itemize}[leftmargin=1.5em, itemsep=2pt]
    \item \textbf{B1}: $\mathbf{P}^{-1}$ exists $\iff \det(\mathbf{P}) \neq 0$, so $k \in \mathbb{R} \setminus \{0, 3\}$ ($k \neq 0$ and $k \neq 3$).
  \end{itemize}
}

\inv{When $k = 0$, what does the transformation represented by $\mathbf{P}$ do to 3D space? What is the dimension of the column space (image) and null space (kernel) when $k = 0$ versus $k = 3$?}

\vspace{10pt}

%── Q5 [REV] ─────────────────────────────────────────────────────────────
\qn{5} \textbf{Inverse Products and Matrix Similarity}
\ans{(a) Proof complete; \quad (b)(i) $\det(\mathbf{M}) = -3$; \quad (ii) $\det(\mathbf{M}^5) = -243$; \quad (iii) $\begin{pmatrix} 81 & 0 \\ 0 & 1 \end{pmatrix}$.}

\method{
  \textbf{(a)}
  \begin{itemize}[leftmargin=1.5em, itemsep=2pt]
    \item \textbf{M1}: Show that $(\mathbf{A}\mathbf{B})(\mathbf{B}^{-1}\mathbf{A}^{-1}) = \mathbf{A}(\mathbf{B}\mathbf{B}^{-1})\mathbf{A}^{-1} = \mathbf{A}\mathbf{I}\mathbf{A}^{-1} = \mathbf{I}$.
    \item \textbf{A1}: State that similarly $(\mathbf{B}^{-1}\mathbf{A}^{-1})(\mathbf{A}\mathbf{B}) = \mathbf{I}$, thus by definition of matrix inverse $(\mathbf{A}\mathbf{B})^{-1} = \mathbf{B}^{-1}\mathbf{A}^{-1}$.
  \end{itemize}
  \textbf{(b)(i)}
  \begin{itemize}[leftmargin=1.5em, itemsep=2pt]
    \item \textbf{M1}: Use determinant product rule: $\det(\mathbf{N}^{-1}\mathbf{M}\mathbf{N}) = \det(\mathbf{N}^{-1})\det(\mathbf{M})\det(\mathbf{N}) = \frac{1}{\det(\mathbf{N})}\det(\mathbf{M})\det(\mathbf{N}) = \det(\mathbf{M})$.
    \item \textbf{A1}: $\det\begin{pmatrix} 3 & 0 \\ 0 & -1 \end{pmatrix} = 3(-1) - 0 = -3 \implies \det(\mathbf{M}) = -3$.
  \end{itemize}
  \textbf{(b)(ii)}
  \begin{itemize}[leftmargin=1.5em, itemsep=2pt]
    \item \textbf{B1}: $\det(\mathbf{M}^5) = (\det\mathbf{M})^5 = (-3)^5 = -243$.
  \end{itemize}
  \textbf{(b)(iii)}
  \begin{itemize}[leftmargin=1.5em, itemsep=2pt]
    \item \textbf{M1}: Observe that $(\mathbf{N}^{-1}\mathbf{M}\mathbf{N})^4 = (\mathbf{N}^{-1}\mathbf{M}\mathbf{N})(\mathbf{N}^{-1}\mathbf{M}\mathbf{N})\cdots = \mathbf{N}^{-1}\mathbf{M}^4\mathbf{N}$.
    \item \textbf{A1}: $\mathbf{N}^{-1}\mathbf{M}^4\mathbf{N} = \begin{pmatrix} 3 & 0 \\ 0 & -1 \end{pmatrix}^4 = \begin{pmatrix} 3^4 & 0 \\ 0 & (-1)^4 \end{pmatrix} = \begin{pmatrix} 81 & 0 \\ 0 & 1 \end{pmatrix}$.
  \end{itemize}
}

\inv{The transformation $\mathbf{N}^{-1}\mathbf{M}\mathbf{N} = \mathbf{D}$ is called diagonalization. The diagonal elements $3$ and $-1$ are the eigenvalues of $\mathbf{M}$. How can this relation be used to compute high matrix powers $\mathbf{M}^k$ rapidly?}

\vspace{10pt}

%── Q6 [REV] ─────────────────────────────────────────────────────────────
\qn{6} \textbf{Consistency Analysis and Parameterised Systems}
\ans{(a) Proof complete; \quad (b)(i) $k = 0$; \quad (ii) $d = 5$; \quad (iii) For $d=5$: sheaf of planes intersecting in a common straight line; for $d \neq 5$: planes form a triangular prism.}

\method{
  \textbf{(a)}
  \begin{itemize}[leftmargin=1.5em, itemsep=2pt]
    \item \textbf{M1}: Determinant expansion: $1(-3 - k) - 1(6 - 3) + 2(2k - (-3))$.
    \item \textbf{A1}: Simplify: $-3 - k - 3 + 4k + 6 = 3k$.
  \end{itemize}
  \textbf{(b)(i)}
  \begin{itemize}[leftmargin=1.5em, itemsep=2pt]
    \item \textbf{B1}: Non-unique solution requires singular matrix $\implies 3k = 0 \implies k = 0$.
  \end{itemize}
  \textbf{(b)(ii)}
  \begin{itemize}[leftmargin=1.5em, itemsep=2pt]
    \item \textbf{M1}: For $k = 0$, examine equations: (1) $x + y + 2z = 4$, (2) $2x - y + z = 1$, (3) $3x + 3z = d$.
    \item \textbf{M1}: Add $(1) + (2)$: $(x+2x) + (y-y) + (2z+z) = 3x + 3z = 4 + 1 = 5$.
    \item \textbf{A1}: Comparing with (3), consistent infinite solutions occur if and only if $d = 5$.
  \end{itemize}
  \textbf{(b)(iii)}
  \begin{itemize}[leftmargin=1.5em, itemsep=2pt]
    \item \textbf{B1}: When $d = 5$: the equations are linearly dependent; since no two normal vectors are parallel, the planes intersect in a single common line (a sheaf of planes).
    \item \textbf{B1}: When $d \neq 5$: $(1)+(2)$ gives $3x + 3z = 5$, contradicting $3x + 3z = d$, so the system is inconsistent. Since no two planes are parallel, they intersect in three mutually parallel lines, forming an open triangular prism.
  \end{itemize}
}

\inv{Explain why testing whether the planes are pairwise parallel is a strictly necessary step before concluding they form a triangular prism. What other geometrical shape could three planes form if two were parallel?}

\newpage
\section*{Section C \quad Proof, Geometry \& Multi-step}

%── Q7 [GEO] ─────────────────────────────────────────────────────────────
\qn{7} \textbf{Inconsistent Systems: The Triangular Prism}
\ans{(a) Inconsistency proof complete; \quad (b) Normals verified non-parallel; \quad (c) Open triangular prism; direction vector $\mathbf{d} = \begin{pmatrix} 1 \\ 3 \\ 5 \end{pmatrix}$.}

\method{
  \textbf{(a)}
  \begin{itemize}[leftmargin=1.5em, itemsep=2pt]
    \item \textbf{M1}: Add equations (1) and (2) to eliminate $z$:
    \[
    (x - 2y + z) + (2x + y - z) = 3x - y = 3 + 4 = 7
    \]
    \item \textbf{A1}: Compare with equation (3): $3x - y = 8$. Subtracting gives $0 = 1$ (contradiction).
    \item \textbf{A1}: Conclude the system is inconsistent, so no points lie on all three planes simultaneously.
  \end{itemize}
  \textbf{(b)}
  \begin{itemize}[leftmargin=1.5em, itemsep=2pt]
    \item \textbf{M1}: Identify normal vectors: $\mathbf{n}_1 = \begin{pmatrix} 1 \\ -2 \\ 1 \end{pmatrix}$, $\mathbf{n}_2 = \begin{pmatrix} 2 \\ 1 \\ -1 \end{pmatrix}$, $\mathbf{n}_3 = \begin{pmatrix} 3 \\ -1 \\ 0 \end{pmatrix}$.
    \item \textbf{A1}: Check scalar multiples: $\mathbf{n}_1 \neq k\mathbf{n}_2$, $\mathbf{n}_1 \neq k\mathbf{n}_3$, $\mathbf{n}_2 \neq k\mathbf{n}_3$. Hence no two planes are parallel.
  \end{itemize}
  \textbf{(c)}
  \begin{itemize}[leftmargin=1.5em, itemsep=2pt]
    \item \textbf{B1}: Since no two planes are parallel and they have no common point, each pair intersects in a line, and the three intersection lines are mutually parallel $\implies$ triangular prism.
    \item \textbf{M1}: Direction of intersection lines is parallel to $\mathbf{n}_1 \times \mathbf{n}_2$:
    \[
    \mathbf{d} = \begin{pmatrix} 1 \\ -2 \\ 1 \end{pmatrix} \times \begin{pmatrix} 2 \\ 1 \\ -1 \end{pmatrix} = \begin{pmatrix} (-2)(-1) - (1)(1) \\ (1)(2) - (1)(-1) \\ (1)(1) - (-2)(2) \end{pmatrix} = \begin{pmatrix} 1 \\ 3 \\ 5 \end{pmatrix}
    \]
    \item \textbf{A1}: $\mathbf{d} = \begin{pmatrix} 1 \\ 3 \\ 5 \end{pmatrix}$ (or any non-zero scalar multiple). Note that $\mathbf{d} \cdot \mathbf{n}_3 = 1(3) + 3(-1) + 5(0) = 0$, confirming parallelism.
  \end{itemize}
}

\inv{What is the cross-sectional area of the triangular prism formed by these three planes? Can you calculate the distance between the parallel lines of intersection?}

\vspace{10pt}

%── Q8 [GEO] ─────────────────────────────────────────────────────────────
\qn{8} \textbf{Dependent Systems: Sheaf of Planes}
\ans{(a) $\det(\mathbf{M}) = 0$; \quad (b) System is consistent ($2\Pi_1 + \Pi_2 = \Pi_3$); normals non-parallel $\implies$ sheaf of planes.}

\method{
  \textbf{(a)}
  \begin{itemize}[leftmargin=1.5em, itemsep=2pt]
    \item \textbf{M1}: Set up matrix and expand determinant:
    \[
    \det(\mathbf{M}) = 1(-1 - 9) - 2(2 - 12) - 1(6 - (-4))
    \]
    \item \textbf{A1}: $-10 - 2(-10) - 10 = -10 + 20 - 10 = 0$.
  \end{itemize}
  \textbf{(b)}
  \begin{itemize}[leftmargin=1.5em, itemsep=2pt]
    \item \textbf{M1}: Form linear combination: consider $2\Pi_1 + \Pi_2$:
    \[
    2(x + 2y - z) + (2x - y + 3z) = 4x + 3y + z
    \]
    \item \textbf{A1}: Check constants: $2(2) + 9 = 13$, matching $\Pi_3$ exactly.
    \item \textbf{B1}: State that the third equation is linearly dependent on the first two.
    \item \textbf{M1}: Check normals: $\mathbf{n}_1 = (1, 2, -1)$ and $\mathbf{n}_2 = (2, -1, 3)$ are not scalar multiples, so $\Pi_1$ and $\Pi_2$ are not parallel and intersect in a line $L$.
    \item \textbf{A1}: Since $2\Pi_1 + \Pi_2 = \Pi_3$, every point satisfying $\Pi_1$ and $\Pi_2$ also satisfies $\Pi_3$. Hence all three planes intersect along the common line $L$ (sheaf of planes).
  \end{itemize}
}

\inv{A sheaf of planes can be visualized as the pages of an open book meeting at the spine. In computer graphics and robotics, how do coordinate systems use pencil/sheaf structures to represent rotational freedom around an axis?}

\newpage

%── Q9 [GEO] ─────────────────────────────────────────────────────────────
\qn{9} \textbf{Line of Intersection Parameterisation}
\ans{(a) $x = 4 - \lambda,\; y = -1 + \lambda$; \quad (b) $\mathbf{r} = \begin{pmatrix} 4 \\ -1 \\ 0 \end{pmatrix} + \lambda \begin{pmatrix} -1 \\ 1 \\ 1 \end{pmatrix}$; \quad (c) Identity verified for all $\lambda$.}

\method{
  \textbf{(a)}
  \begin{itemize}[leftmargin=1.5em, itemsep=2pt]
    \item \textbf{M1}: Substitute $z = \lambda$ into equations (1) and (2):
    \[
    x + 2y = 2 + \lambda \quad \text{and} \quad 2x - y = 9 - 3\lambda
    \]
    \item \textbf{M1}: Multiply second equation by 2 and add:
    \[
    (x + 2y) + 2(2x - y) = 5x = (2 + \lambda) + (18 - 6\lambda) = 20 - 5\lambda
    \]
    \item \textbf{A1}: $x = 4 - \lambda$.
    \item \textbf{A1}: Substitute into $y = 2x - 9 + 3\lambda = 2(4 - \lambda) - 9 + 3\lambda = 8 - 2\lambda - 9 + 3\lambda = -1 + \lambda$.
  \end{itemize}
  \textbf{(b)}
  \begin{itemize}[leftmargin=1.5em, itemsep=2pt]
    \item \textbf{M1}: Write in vector form $\begin{pmatrix} x \\ y \\ z \end{pmatrix} = \begin{pmatrix} 4 - \lambda \\ -1 + \lambda \\ \lambda \end{pmatrix}$.
    \item \textbf{A1}: Split into position and direction vectors:
    \[
    \mathbf{r} = \begin{pmatrix} 4 \\ -1 \\ 0 \end{pmatrix} + \lambda \begin{pmatrix} -1 \\ 1 \\ 1 \end{pmatrix}
    \]
  \end{itemize}
  \textbf{(c)}
  \begin{itemize}[leftmargin=1.5em, itemsep=2pt]
    \item \textbf{M1}: Substitute $x = 4 - \lambda$, $y = -1 + \lambda$, $z = \lambda$ into LHS of $\Pi_3$:
    \[
    4(4 - \lambda) + 3(-1 + \lambda) + (\lambda)
    \]
    \item \textbf{A1}: $16 - 4\lambda - 3 + 3\lambda + \lambda = 13 \equiv 13$ (RHS). The equation holds for all $\lambda \in \mathbb{R}$.
  \end{itemize}
}

\inv{Verify that the direction vector $\mathbf{d} = (-1, 1, 1)^T$ is perpendicular to all three normal vectors $\mathbf{n}_1, \mathbf{n}_2, \mathbf{n}_3$ by evaluating each scalar product $\mathbf{d} \cdot \mathbf{n}_i$. Why must this always hold?}

\vspace{10pt}

%── Q10 [GEO] ─────────────────────────────────────────────────────────────
\qn{10} \textbf{Parallel Planes and Transversal Intersections}
\ans{(a) $\mathbf{n}_2 = 2\mathbf{n}_1$ and $15 \neq 8 \implies$ parallel and distinct; \quad (b) Common direction vector $\mathbf{d} = \begin{pmatrix} -1 \\ 1 \\ 1 \end{pmatrix}$; \quad (c) Two parallel intersection lines; zero solutions (inconsistent).}

\method{
  \textbf{(a)}
  \begin{itemize}[leftmargin=1.5em, itemsep=2pt]
    \item \textbf{M1}: Normal vectors are $\mathbf{n}_1 = \begin{pmatrix} 2 \\ -1 \\ 3 \end{pmatrix}$ and $\mathbf{n}_2 = \begin{pmatrix} 4 \\ -2 \\ 6 \end{pmatrix} = 2\mathbf{n}_1$.
    \item \textbf{A1}: Since $\mathbf{n}_2$ is a scalar multiple of $\mathbf{n}_1$, the planes are parallel. Comparing constants: $\Pi_2$ can be written $2x - y + 3z = 7.5 \neq 4$, so the planes are distinct and never intersect.
  \end{itemize}
  \textbf{(b)}
  \begin{itemize}[leftmargin=1.5em, itemsep=2pt]
    \item \textbf{M1}: The direction vector of $L_1 = \Pi_1 \cap \Pi_3$ is parallel to $\mathbf{n}_1 \times \mathbf{n}_3$:
    \[
    \mathbf{d} = \begin{pmatrix} 2 \\ -1 \\ 3 \end{pmatrix} \times \begin{pmatrix} 1 \\ 2 \\ -1 \end{pmatrix} = \begin{pmatrix} 1 - 6 \\ 3 - (-2) \\ 4 - (-1) \end{pmatrix} = \begin{pmatrix} -5 \\ 5 \\ 5 \end{pmatrix} \sim \begin{pmatrix} -1 \\ 1 \\ 1 \end{pmatrix}
    \]
    \item \textbf{A1}: State direction vector $\mathbf{d} = \begin{pmatrix} -1 \\ 1 \\ 1 \end{pmatrix}$ (or scalar multiple).
    \item \textbf{M1}: For $L_2 = \Pi_2 \cap \Pi_3$, the direction vector is $\mathbf{n}_2 \times \mathbf{n}_3 = 2(\mathbf{n}_1 \times \mathbf{n}_3) = 2\mathbf{d}$.
    \item \textbf{A1}: Since both intersection lines share the direction $\mathbf{d}$ and lie in distinct parallel planes, $L_1$ and $L_2$ are parallel lines.
  \end{itemize}
  \textbf{(c)}
  \begin{itemize}[leftmargin=1.5em, itemsep=2pt]
    \item \textbf{B1}: Geometric arrangement: Two parallel planes cut transversally by a third plane, resulting in two parallel lines of intersection.
    \item \textbf{B1}: Because $\Pi_1$ and $\Pi_2$ do not intersect, there are 0 solutions common to all three equations (the system is inconsistent).
  \end{itemize}
}

\inv{What happens if $\Pi_2$ is replaced by $4x - 2y + 6z = 8$? In that case, what is the geometric relationship between $\Pi_1$ and $\Pi_2$, and how many solutions does the three-plane system have?}

\end{document}
